\section{When counting errors imply mixing}
\label{sec:countmixing}

Mixing means that observing a system now and much later becomes
asymptotically independent. For a probability-preserving flow,
\[
 \operatorname{Corr}_t(f,h)=
 \int f\,(h\circ g_t)\,\dd\mu
       -\left(\int f\,\dd\mu\right)\left(\int h\,\dd\mu\right).
\]
Mixing asks that this tend to zero for appropriate observables;
a rate bounds how fast it does so in terms of stated norms.
A scalar count need not retain both observations. The examples
below identify counts that do, and show what is lost by counts
that forget this information.

A counting remainder can imply mixing when the counts determine an
appropriate operator, or when they retain sufficiently uniform information
about both endpoints of orbit segments. The following elementary
converses make these hypotheses and their quantitative losses explicit.
They do not identify the simple-curve counting function of
Section~\ref{sec:quantitativecounting} with such an operator.

\subsection{A direct converse in the compact Selberg setting}
\label{subsec:countmixselberg}

There is a particularly close answer to the original trace-formula
question. Let $X=\Gamma\backslash\mathbb H^2$ be a closed connected
oriented surface of curvature $-1$. Let $\pi_X(L)$ count primitive
hyperbolic conjugacy classes in $\Gamma$ of length at most $L$.
Thus the two orientations are counted separately: $[\gamma]$ and
$[\gamma^{-1}]$ are distinct. This is the primitive-class convention
in the Selberg product; see \cite[Section~2.1]{PollicottSharpZeta}.

\begin{proposition}[A prime-geodesic remainder forces a mixing rate]
\label{prop:countmixselberg}
Suppose that, for some $0<\delta<\tfrac12$,
\begin{equation}
 \pi_X(L)=\operatorname{Li}(e^L)+O(e^{(1-\delta)L}).
 \label{eq:countmixselbergcount}
\end{equation}
Then the first positive Laplace eigenvalue satisfies
\begin{equation}
 \lambda_1(X)\geq\delta(1-\delta).
 \label{eq:countmixselberggap}
\end{equation}
For every $0<r<\delta$, the unit-speed geodesic flow $g_t$ on $T^1X$
therefore satisfies
\begin{equation}
 |\operatorname{Corr}_t(f,h)|
 \leq C_{X,r}S_m(f)S_m(h)e^{-rt}
 \label{eq:countmixselbergmixing}
\end{equation}
for some finite Sobolev order $m$. Here $S_m$ is an $L^2$ Sobolev norm
and the correlation uses normalized Liouville measure.
\end{proposition}

\begin{proof}
The proof uses a spectral dictionary: a small positive Laplace
eigenvalue would create an additional pole in a transformed
geodesic count. A sufficiently small counting remainder rules out
that pole. An external matrix-coefficient theorem then converts
the eigenvalue gap to decay of correlations. Here
$\operatorname{Li}(x)=\int_2^x\dd u/\log u$ (the lower cutoff
changes only a constant), and a Sobolev norm measures the $L^2$
sizes of a function and its derivatives up to a fixed order.

Put $\alpha=1-\delta>\tfrac12$. For $\Re s>1$, define
\[
 F(s)=\sum_{P\ \mathrm{primitive}}\ell(P)e^{-s\ell(P)}
     =\int_0^\infty t e^{-st}\,d\pi_X(t).
\]
Above a fixed positive cutoff $L_0$, write
$\pi_X(t)=\operatorname{Li}(e^t)+E(t)$ with
$E(t)=O(e^{\alpha t})$. Since
$d\operatorname{Li}(e^t)=e^t\,dt/t$, the main term contributes
$e^{-(s-1)L_0}/(s-1)$. Integration by parts gives the error contribution,
up to an entire cutoff term, as
\[
 \int_{L_0}^\infty (st-1)e^{-st}E(t)\,dt.
\]
This integral is locally uniformly convergent for $\Re s>\alpha$.
Consequently $F(s)-1/(s-1)$ is holomorphic there.

The logarithmic derivative of the Selberg product is
\begin{equation}
 \frac{Z_X'(s)}{Z_X(s)}
 =\sum_P\sum_{k\geq1}
       \frac{\ell(P)e^{-sk\ell(P)}}{1-e^{-k\ell(P)}}.
 \label{eq:countmixselberglogderivative}
\end{equation}
Subtracting $F(s)$ leaves a function holomorphic for $\Re s>\tfrac12$:
the terms with $k\geq2$ have exponent at least $2\Re s>1$, and the
$k=1$ denominator correction supplies an extra factor $e^{-\ell(P)}$.
The positive systole and the counting bound justify local uniform
convergence. Hence $Z_X'/Z_X-1/(s-1)$ is holomorphic for
$\Re s>\alpha$, excluding zeros of $Z_X$ there except at $1$.

Every Laplace eigenvalue $\lambda_j=s_j(1-s_j)<\tfrac14$, with
$\tfrac12<s_j<1$, gives a zero $s_j$ of $Z_X$; see
\cite[Theorems~1.1 and~4.1]{BJPZeta}. Thus $s_j\leq1-\delta$,
which proves \eqref{eq:countmixselberggap}; equality is allowed.
Set $s_*=\tfrac12$ if there is no such positive eigenvalue, and otherwise
let $s_*$ be the largest such parameter. The matrix-coefficient estimate
of \cite[Theorem~1.1]{EdwardsOhMixing}, specialized to a compact surface,
gives Sobolev correlation decay at every rate below $1-s_*$.
Since $s_*\leq1-\delta$, this proves
\eqref{eq:countmixselbergmixing}. The time normalization is
$a_t=\operatorname{diag}(e^{t/2},e^{-t/2})$ in
$\mathrm{PSL}_2(\mathbb R)$.
\end{proof}

Two qualifications delimit this converse. First,
\eqref{eq:countmixselbergcount} contains only the leading term:
an error estimate \emph{after subtracting} exceptional contributions
$\operatorname{Li}(e^{s_jL})$ does not exclude those small eigenvalues.
Second, the conclusion concerns the complete primitive spectrum.
Applying it to simple curves or coefficient-filtered counts would require
a corresponding spectral identity. If geometric geodesics are counted
without orientation, the main term is
$\tfrac12\operatorname{Li}(e^L)$; doubling that count recovers the stated
convention. The exponent is quantitative, while explicit numerical or
family-uniform mixing constants require additional estimates.

\subsection{A quantitative converse in homogeneous dynamics}

The next paragraph is an optional comparison with an arithmetic
setting. An averaging operator translates a function by group
elements in $B$ and averages the results using probability measure
$\beta_B$. Its norm on mean-zero functions measures how much
nonconstant information survives that average. Discrepancy is the
difference between a normalized count and its predicted volume.
The cited theorem bounds this operator norm by a localized $L^2$
discrepancy. Its $S$-arithmetic hypotheses, involving real and
selected non-Archimedean factors, are external prerequisites;
the finite-state and localized arguments below do not require them.

There is an established arithmetic converse. In the $S$-arithmetic
setting of Gorodnik--Nevo~\cite[Theorems 1.4 and 5.2]{GNDiscrepancy},
let $E_{S,m}(r,B)$ denote the $L^2$ discrepancy of normalized rational
point counts in translated real balls of radius $r$ and an
$S$-adic set $B$. For a discretely embedded irreducible representation
$\pi$ in the mean-zero automorphic spectrum, their estimate is
\begin{equation}
 \|\pi(\beta_B)\|
 \leq 2[\Gamma_S:\Gamma_{S,m}]
       m_{G_\infty}(B(e,r))^{-1}E_{S,m}(r,B),
 \qquad 0<r\leq r_0(\pi).
 \label{eq:countmixGN}
\end{equation}
The ambient group and congruence quotients satisfy the standing
algebraic hypotheses of that paper, including a connected, simply
connected $\mathbb Q$-simple group in the notation displayed here.
Here the real factor is noncompact, the places in $S$ are unramified,
$B$ is bounded, positive-measure, and bi-invariant under the specified
hyperspecial compact subgroup, and $\beta_B$ is its normalized Haar
measure. The small-radius threshold depends on $\pi$.
Thus sufficiently strong localized discrepancy controls averaging
operators quantitatively. Their application to forms of $\mathrm{SL}_2$
uses arithmetic counting to obtain a spectral gap. This statement is
not a general implication from one scalar count to mixing of a flow.


\subsection{A converse for a finite Markov model}
\label{subsec:countmixfinite}

Let $A$ be an irreducible nonnegative integer $q\times q$ adjacency
matrix, allowing parallel edges, and let $\lambda>0$ be its Perron
eigenvalue. The number $F_n$ of periodic edge sequences fixed by the
$n$th shift, with a distinguished starting position, is
\[
 F_n=\tr(A^n).
\]
It includes repetitions and is not the number of primitive orbits.
The associated identity
\[
 \zeta_A(z)=
 \exp\left(\sum_{n\geq1}\frac{\tr(A^n)}n z^n\right)
 =\frac1{\det(I-zA)}
\]
is the usual finite-state trace formula; see
\cite[Sections 1 and 4]{SharpZeta}.

\begin{proposition}[A trace-counting remainder forces a gap]
\label{prop:countmixfinite}
Suppose that, for a specified $\kappa>0$,
\begin{equation}
 \tr(A^n)=\lambda^n+
          O\bigl((\lambda e^{-\kappa})^n\bigr).
 \label{eq:countmixtrace}
\end{equation}
Then every eigenvalue of $A$ other than its simple Perron eigenvalue
satisfies
\begin{equation}
 |\lambda_j|\leq\lambda e^{-\kappa}.
 \label{eq:countmixgap}
\end{equation}
\end{proposition}

\begin{proof}
Put $\rho=\lambda e^{-\kappa}$. The power series
\[
 E(z)=\sum_{n\geq1}\bigl(\tr(A^n)-\lambda^n\bigr)z^n
\]
is holomorphic for $|z|<\rho^{-1}$. On its initial convergence disk,
\[
 E(z)=\sum_{j=2}^q\frac{z\lambda_j}{1-z\lambda_j},
\]
where algebraic multiplicities are included. A nonzero $\lambda_j$
produces a pole at $z=\lambda_j^{-1}$. Distinct eigenvalues produce
different poles, and repeated copies of one eigenvalue have the same
nonzero residue and add. No pole can therefore lie inside
$|z|<\rho^{-1}$. Zero eigenvalues already satisfy the conclusion.
\end{proof}

For a complete two-state example, take
$A=\left(\begin{smallmatrix}2&1\\1&2\end{smallmatrix}\right)$.
Its eigenvalues are $3$ and $1$, so $\tr A^n=3^n+1$ and
$\kappa=\log3$ is a valid saving. The transition matrix is
$P=A/3$, with stationary probabilities $(1/2,1/2)$.
The mean-zero observable $f=(1,-1)$ satisfies $Pf=f/3$,
so its stationary correlation at lag $n$ is exactly $3^{-n}$.
This makes the connection between the secondary eigenvalue,
counting error and loss of memory explicit.

To convert the gap into a mixing statement, choose positive Perron
vectors $r,l$ normalized by $\sum_i l_ir_i=1$, and set
\[
 P_{ij}=\frac{A_{ij}r_j}{\lambda r_i},
 \qquad \pi_i=l_ir_i,\qquad \Pi=\mathbf1\pi^T.
\]
Then $P$ is stochastic, $\pi$ is stationary, and $P$ is similar to
$A/\lambda$. For every $0<\kappa'<\kappa$,
\begin{equation}
 \|P^n-\Pi\|_\infty\leq C_{\kappa',A}e^{-\kappa'n}.
 \label{eq:countmixparry}
\end{equation}
At the endpoint $\kappa$ a polynomial factor can arise from Jordan
blocks. In particular, the stationary Markov chain satisfies
\[
 \left|\mathbb E_\pi[f(X_0)g(X_n)]
       -\mathbb E_\pi f\,\mathbb E_\pi g\right|
 \leq C_{\kappa',A}e^{-\kappa'n}\|f\|_\infty\|g\|_\infty.
\]
Fixed-width cylinder observables have the same decay after the
corresponding time offset. Approximation by cylinder functions gives
mixing of the shift, and quantitative approximation gives a rate for
H\"older observables with an explicit regularity loss. For example, on
the two-sided stationary vertex shift use the metric $d_\theta(x,y)
=\theta^{N(x,y)}$, where $0<\theta<1$ and $N(x,y)$ is the least
$|j|$ for which $x_j\ne y_j$. A $\beta$-H\"older test is approximated
by conditioning on coordinates $-p,\ldots,p$ with error
$O(\theta^{\beta p})$. The Markov property bounds the covariance of
the two resulting cylinder tests at lag $n>2p$ by
$O(e^{-\kappa'(n-2p)})$ times their sup norms. Taking
$p=\lfloor\kappa'n/(2\kappa'+\beta|\log\theta|)\rfloor$ gives the rate
\[
 \exp\left(-\frac{\kappa'\beta|\log\theta|}
 {2\kappa'+\beta|\log\theta|}\,n\right).
\]
Parallel-edge observations are handled by conditioning on the adjoining
vertices; a bounded change in the time offset changes only the constant.

The matrix constant can itself be bounded explicitly. If
$r_0=e^{-\kappa}<r<1$, the resolvent integral around $|z|=r$ gives
\begin{equation}
 \|P^n-\Pi\|_\infty
 \leq
 \frac{r\,q!(1+r)^{q-1}}
      {(1-r)(r-r_0)^{q-1}}\,r^n.
 \label{eq:countmixmatrixconstant}
\end{equation}
Indeed, on this circle the determinant has modulus at least
$(1-r)(r-r_0)^{q-1}$, while the row-sum norm of its adjugate is
at most $q!(1+r)^{q-1}$. Thus a numerical counting saving produces
a numerical mixing estimate in this model.

\subsection{Which scalar count is being estimated?}
\label{subsec:countmixscalar}

An exact count of unrestricted paths need not imply mixing.
For example, with $m\geq2$,
\[
 A=\begin{pmatrix}0&m\\m&0\end{pmatrix},
 \qquad \mathbf1^T A^n\mathbf1=2m^n.
\]
This count has no error at all. Its normalized transition chain
alternates between two states; the correlation of $(1,-1)$ with
itself is $(-1)^n$. The closed-path count detects what the total
path count misses:
\[
 \tr(A^n)=m^n+(-m)^n.
\]
Thus a scalar sequence can suffice when it is a complete trace
encoding of a suitable operator. The decisive issue is what the
chosen count retains.

Continuous time introduces a separate return-time issue. The full
shift on $m$ symbols has exactly $m^n$ fixed points for its $n$th
iterate and is mixing for its Bernoulli measure. Suspend it with
constant roof $1$:
\[
 Y=(\Sigma\times[0,1])/
       \bigl((x,1)\sim(\sigma x,0)\bigr).
\]
For the suspension flow $\phi^t$, the mean-zero function
$F([x,u])=e^{2\pi i u}$ is well defined and satisfies
\[
 F\circ\phi^t=e^{2\pi it}F,\qquad
 \int\overline F\,F\circ\phi^t\,\dd\nu=e^{2\pi it}.
\]
The flow is therefore not mixing. Its physical periods lie in a
lattice, so this is not an example of an unoscillating continuous-time
prime-orbit asymptotic. It shows why a discrete-time counting converse
does not automatically pass to a suspension. Ruling out arithmetic
return times is necessary; quantitative flow estimates also require
control of approximate eigenfunctions or high-frequency return-time
operators. The role of these additional geometric conditions is
discussed in~\cite{DolgopyatMixing}.

\subsection{A qualitative converse from full geodesic counting}

There is also a simple geometric deduction. In
Babillot~\cite[Theorem~1]{BabillotMixing},
the universal cover is a simply connected complete Riemannian manifold
with sectional curvature bounded above by a negative constant, and
the acting discrete isometry group is non-elementary. A quasi-product
measure means a flow-invariant measure whose geodesic current is in the
measure class of a product of its two boundary measures. For such a
measure, nonarithmeticity of the closed length spectrum implies mixing.
Assume a nonzero finite such measure, normalized to probability,
and suppose the full
primitive closed-orbit count is finite at each cutoff and satisfies
\[
 \pi(T)\sim C\,\frac{e^{hT}}{T},\qquad C,h>0,
\]
as $T\longrightarrow\infty$ through all real values. If all periods
belonged to $a\ZZ$, $a>0$, then
\[
 \pi(na+a/4)=\pi(na+3a/4),
\]
whereas the ratio of their proposed main terms tends to $e^{ha/2}>1$.
This contradiction proves nonarithmeticity and hence qualitative
mixing under the stated structural hypotheses. No rate follows from
this argument. For a polynomial simple-curve main term $CT^d$,
the corresponding fixed-shift ratio tends to one, so this argument
does not transfer to simple-curve counting.

\subsection{Localized counts and correlations}
\label{subsec:countmixlocalized}

A useful general converse applies once geometric counts are related
to joint transition masses. Let $(Y,d,\mu)$ be a compact metric
probability space, with a measure-preserving flow $\phi^t$. Suppose
that for every $0<h\leq1$ there is a measurable partition
$\mathcal P_h=\{P_1,\ldots,P_{N(h)}\}$ with
\[
 \operatorname{diam}(P_i)\leq h,\qquad N(h)\leq B h^{-D}.
\]
Zero-measure cells may be discarded. Define
$p_{ij}(t,h)=\mu(P_i\cap\phi^{-t}P_j)$.

\begin{proposition}[A quantitative localized converse]
\label{prop:countmixlocalized}
Assume that geometric weighted counts give, uniformly in $i,j,h,t$,
\begin{equation}
 |p_{ij}(t,h)-\mu(P_i)\mu(P_j)|
       \leq A h^{-m}e^{-\kappa t},
 \qquad m\geq0,\quad\kappa>0.
 \label{eq:countmixentry}
\end{equation}
For $0<\beta\leq1$ and bounded $\beta$-H\"older functions $f,g$,
\begin{equation}
 \left|\int f\,g\circ\phi^t\,\dd\mu-\int f\,\dd\mu\int g\,\dd\mu\right|
 \leq C\|f\|_{C^\beta}\|g\|_{C^\beta}e^{-\eta t},
 \qquad
 \eta=\frac{\kappa\beta}{m+2D+\beta}.
 \label{eq:countmixrate}
\end{equation}
\end{proposition}

\begin{proof}
The proof balances spatial resolution against counting error.
A $\beta$-H\"older function satisfies
$|f(x)-f(y)|\le C d(x,y)^\beta$, so replacing it by a constant
on a cell of diameter $h$ costs $O(h^\beta)$. Smaller cells
improve that approximation but create more ordered cell pairs,
which amplify an entrywise counting error. The choice of $h$
at the end makes these two effects comparable.

Replace $f,g$ by their cellwise averages $f_h,g_h$. Their uniform
approximation errors are at most $[f]_\beta h^\beta$ and
$[g]_\beta h^\beta$. Invariance bounds the resulting correlation
error by $C\|f\|_{C^\beta}\|g\|_{C^\beta}h^\beta$.
For the step functions the centered correlation is
\[
 \sum_{i,j}f_i g_j
       \bigl(p_{ij}(t,h)-\mu(P_i)\mu(P_j)\bigr).
\]
There are at most $B^2h^{-2D}$ terms. Equation~\eqref{eq:countmixentry}
therefore bounds the original correlation by
\[
 C\|f\|_{C^\beta}\|g\|_{C^\beta}
       \bigl(h^\beta+h^{-(m+2D)}e^{-\kappa t}\bigr).
\]
Choose $h=e^{-\kappa t/(m+2D+\beta)}$.
\end{proof}

If the right side of~\eqref{eq:countmixentry} contains the additional
factor $\mu(P_i)\mu(P_j)$, its sum costs one, and the exponent
improves to $\kappa\beta/(m+\beta)$. The normalization of a localized
count therefore materially affects the resulting rate.

The substantive hypothesis is the geometric identification of counts
with $p_{ij}$. In an application one needs a formula of the form
\[
 p_{ij}(t,h)=
 \frac{\text{weighted segment count from }P_i\text{ to }P_j}
      {\text{normalizing growth}}
 +\text{controlled error}.
\]
Jacobian weights, multiplicities, boundary effects, and time-window
errors must be controlled uniformly at both endpoints and at shrinking
resolution $h$. If the comparison error satisfies the same bound as
the normalized counting error, their constants add and the proposition
applies. The proposition does not supply this comparison for an
arbitrary raw counting problem.

\subsection{Nonescape and the counting time parameter}
\label{subsec:countmixnoncompact}

There is an explicit version with a noncompact tail. Suppose an
exhaustion $K_H$ satisfies
$\mu(Y\setminus K_H)\leq C e^{-aH}$, with $a>0$.
Assume that $K_H$ has at most $B e^{bH}h^{-D}$ cells of diameter
$h$, and that the entrywise bound is
$A e^{cH}h^{-m}e^{-\kappa t}$, where $b,c,m\geq0$.
For bounded globally $\beta$-H\"older observables the preceding
argument, extending the step approximations by zero off $K_H$, gives
\[
 |\operatorname{Corr}_t(f,g)|
 \leq C\|f\|_{C^\beta}\|g\|_{C^\beta}
 \left(h^\beta+e^{-aH}
       +e^{(2b+c)H}h^{-(m+2D)}e^{-\kappa t}\right).
\]
The tail error uses invariance and the $L^1$ approximation bound;
no uniform approximation on the discarded region is required.
Taking $h=e^{-\eta t/\beta}$ and $H=\eta t/a$ yields
\begin{equation}
 \boxed{\displaystyle
 \eta=
 \frac{\kappa}
 {1+(2b+c)/a+(m+2D)/\beta}.}
 \label{eq:countmixcusprate}
\end{equation}
These are declared sufficient hypotheses, not estimates established
here for Teichm\"uller flow.

Finally, under the counting-time conversion $R=e^t$, a relative power
saving becomes an exponential one:
\[
 \frac{O(R^{d-\delta})}{R^d}=O(e^{-\delta t}).
\]
This explains why the exponents in effective counting and mixing
are related. The conversion alone gives no transition-mass identity.
A converse for our curve-counting program would require the estimates
to retain independent endpoint or sector weights, with controlled
dependence on localization and on excursions into the noncompact
part of the parameter space. Alternatively, an appropriate trace or
transfer-operator representation could provide a spectral converse.
Neither representation follows merely from a good remainder for the
single unweighted simple-length count.

\clearpage
